\section{Conclusion}

\subsection{Conclusion on Simulated Annealing}

The Simulated Annealing algorithm demonstrated robust performance across the eight cooling tower optimisation scenarios, consistently satisfying the volume constraint with errors below 0.5\,m$^3$ in all cases and confirming the effectiveness of the penalty-based approach.

In terms of solution quality, Case~8 (Hyperbolic Fit) achieved the minimum surface area of 5\,765.62\,m$^2$, representing a 26\% reduction compared to the reference case and highlighting the advantage of parametric formulations that reduce the search space dimension while guaranteeing physically realistic geometries. The convergence behaviour across all cases exhibited characteristic staircase patterns, with most improvements occurring within the first 2\,000 iterations; the algorithm successfully escaped local minima in early stages due to the high initial temperature. Sensitivity to initial conditions was observed in Case~5 (Cylindrical Start), where the deliberately poor starting point led to a 2.3\% increase in area compared to the reference case, indicating that SA's final solution retains some dependence on the initialisation despite its stochastic exploration mechanism. The most significant limitation emerged in Case~6 (High Volume), where the absence of explicit geometric constraints---monotonicity, convexity---allowed the SA to produce unrealistic zigzag profiles. This confirms that penalty functions alone are insufficient for shape optimisation problems requiring smooth geometries.

In summary, SA is well-suited for cooling tower optimisation when combined with appropriate geometric constraints or parametric formulations. The computational cost remains low (1.6--6.2\,s per case), making it practical for engineering design applications.

\subsection{Conclusion on Particle Swarm Optimisation}

The PSO algorithm exhibited strong exploration capabilities and reliable constraint satisfaction across all eight cases under the penalty-based formulation used in this study. Volume errors remained below 1.2\,m$^3$ in all cases, confirming that the penalty weight $\rho = 10^5$ is sufficient to guide the swarm toward the feasible manifold. However, PSO consistently produced larger surface areas than BFGS, indicating that while the swarm successfully satisfies the volume constraint, it lacks the local refinement precision required to minimise the objective function effectively.

The best performance was achieved on Case~8 (Hyperbolic Fit), which produced the lowest surface area of all algorithm--case combinations (5\,742.55\,m$^2$) with near-zero volume error, confirming that the low-dimensional parametric space enables effective swarm convergence. The two-phase convergence pattern---rapid descent followed by plateau---was consistently observed, with Cases~2, 6, and~8 converging within 30--50 iterations while Cases~1, 3, and~5 required over 500 iterations. The computational cost constitutes the most significant limitation of PSO as deployed in this study: PSO required on average 133\,800 function evaluations per case---$14\times$ more than BFGS and $8\times$ more than SA---yet produced larger surface areas than BFGS in seven out of eight cases, yielding a poor cost-effectiveness ratio that limits PSO's suitability for problems where function evaluations are expensive.

In summary, PSO's global exploration capability ensures reliable constraint satisfaction but proves insufficient to compete with gradient-based methods in terms of solution quality and computational efficiency. Future work should investigate local refinement strategies, adaptive inertia schedules, or hybrid PSO--BFGS approaches to address the precision gap observed in this study.

\subsection{Conclusion on BFGS}

The BFGS quasi-Newton method emerged as the most effective algorithm in this study, achieving the best or near-best surface area in seven out of eight cases with perfect volume constraint satisfaction (zero error in all cases) and the lowest computational cost.

Its peak performance was observed on Case~8 (Hyperbolic Fit), where BFGS converged in just 12 iterations (0.04\,s) to an area of 5\,961.77\,m$^2$ with zero volume error---the most efficient configuration in the entire study. Across all well-conditioned cases, BFGS exhibited superlinear convergence, with typical convergence achieved within 20--30 iterations; the gradient norm provides a natural, adaptive stopping criterion, in contrast to SA, which relies on a fixed iteration budget, and PSO, which terminates adaptively based on swarm stagnation. The geometric quality of the BFGS solutions was consistently superior: in Case~6 (High Volume), BFGS achieved a 4.7$\times$ lower surface area than SA, attributable entirely to the inherent geometric regularity of gradient-based search. The principal limitation was observed in Case~3 (Full Freedom, $D = 19$), where the method's vulnerability to ill-conditioning and local trapping in high-dimensional spaces led to convergence at a physically unrealistic local minimum. For such problems, BFGS should be combined with multi-start strategies or hybridised with a global search phase.

In summary, BFGS is the recommended algorithm for cooling tower optimisation when gradient information is available and the search space is well conditioned. Its deterministic nature, adaptive termination, and inherent geometric smoothness make it ideally suited for structural engineering applications where solution quality, reproducibility, and computational efficiency are paramount.

\subsection{Overall Conclusion}

This study benchmarked three optimisation algorithms---BFGS, Simulated Annealing, and Particle Swarm Optimisation---on the structural surface area minimisation of a discretised hyperboloidal cooling tower across eight study cases of increasing complexity. The investigation reveals a clear performance hierarchy: BFGS $>$ SA $\gg$ PSO, with the following overarching conclusions:

\begin{enumerate}
    \item \textbf{Algorithm selection matters, but problem formulation matters more.} Case~8 (Hyperbolic Parametrisation, $D = 3$) yielded the best result for all three methods, demonstrating that embedding physical knowledge into the parametrisation is more effective than algorithmic sophistication.
    
    \item \textbf{Gradient-based methods are naturally suited to shape optimisation.} BFGS's local search produces smooth, structurally plausible profiles without requiring explicit geometric penalties, whereas stochastic methods can generate unrealistic geometries unless additional constraints are imposed.
    
    \item \textbf{Penalty-based constraint handling is effective across all three methods.} The static penalty weight $\rho = 10^5$ is sufficient to guide all three algorithms toward the feasible manifold, with volume errors below 1.2\,m$^3$ for PSO, below 0.75\,m$^3$ for SA, and zero for BFGS in seven out of eight cases. However, PSO's distributed evaluation strategy limits its ability to exploit local curvature near the feasible manifold, suggesting that a hybrid penalty--refinement approach would improve surface area precision.
    
    \item \textbf{Computational efficiency spans three orders of magnitude.} BFGS achieves convergence in $\mathcal{O}(10^2)$ function evaluations, SA in $\mathcal{O}(10^4)$, and PSO in $\mathcal{O}(10^5)$, with BFGS delivering superior solutions at the lowest cost in the majority of cases.
\end{enumerate}

\subsection{Future Work}

Building on the limitations identified in Section~5.4.4, several directions for future investigation are proposed:

\begin{itemize}
    \item \textbf{Constraint handling}: As noted in the limitations discussion, the penalty weight $\rho = 10^5$ is suboptimal for the stochastic methods. Future work should investigate augmented Lagrangian methods, feasibility-preserving operators for PSO, and adaptive penalty schedules to improve performance on equality-constrained problems.
    
    \item \textbf{Hybrid algorithms}: Combine global exploration (SA or PSO) with local refinement (BFGS) in a two-phase strategy, using the stochastic method to identify promising basins and BFGS to converge within them.
    
    \item \textbf{Structural constraints}: Extending the cost function to incorporate stress, buckling, and wind-load criteria---identified in Section~5.4.4 as a key omission of the present study---would transform the problem into a multi-objective optimisation that more closely reflects real engineering requirements.
    
    \item \textbf{Higher-fidelity models}: Replace the conical frustum discretisation with finite element shell models to capture bending stiffness, membrane stress distributions, and nonlinear geometric effects.
    
    \item \textbf{Parametric exploration}: Investigate alternative parametric families beyond the hyperboloid of one sheet (e.g., polynomial profiles, B-spline representations) to assess whether further surface area reductions are achievable while maintaining structural integrity.
\end{itemize}
